\section{Supplementary material}
\label{sec:appendix}

\subsection{Post-registration decisions}
\label{sec:appendix-deviations}

Every decision taken after the preregistration lock is logged with its date and reason in the project's decision log; Table~\ref{tab:deviations} lists those that bear on the reported analyses. None changes a hypothesis, a threshold, a sample size or the evaluation of a confirmatory contrast. Two are readings of the protocol where its text left a detail open (the scope of the saved propagation diagnostics and the definition of a secondary metric), fixed in writing before the code freeze. Two concern an implementation defect that stopped the dry run: a directory holding an unreleased memory-mapped file could not be renamed on the operating system used; it was diagnosed, fixed and re-tagged before any model ran on TEST, and the dry run was repeated from scratch. The last is the decision not to run the conditional causal study (Section~\ref{sec:discussion-limitations}).

\begin{table}[ht]
\centering
\caption{Post-registration decisions that bear on the reported analyses (identifiers of the project's decision log).}
\label{tab:deviations}
\small
\begin{tabular}{l p{12cm}}
\toprule
Entry & Decision (DECISIONS.md) \\
\midrule
D78 & Plan review reading of plan section 7.4: full state tensors are not saved; cumulative-state distances and intermediate decodings exist for the primary configurations only \\
D79 & Implementation reading of the naked-single F1 (secondary metric), fixed before the code freeze \\
D81 & Post-freeze fix before any TEST run: retry of a directory rename on Windows \\
D82 & Correction of D81: the rename failed because of an unreleased memory map; fixed at the source and re-tagged before any TEST run \\
D84 & E4 dropped: the plan section 9.1 entry condition is not met given the TEST readout and equivariance results; SELECT and E4-TEST were never generated \\
\bottomrule
\end{tabular}

\end{table}

\subsection{Reproducibility}
\label{sec:appendix-repro}

The protocol is registered at OSF (\url{https://osf.io/m28fx}) and frozen in the repository under the tag \texttt{prereg-v3.7}; the analysis code that produced every TEST output is frozen under the tag \texttt{wp2-code}. Each confirmatory run records its code commit, environment and input hashes, and the evaluation refused to start unless every run record was valid at the tag commit. All tables, figures and numbers in this paper are generated by one script from the committed result files, and a check verifies that the generated files are current.

\subsection{Full result tables}
\label{sec:appendix-tables}

Table~\ref{tab:e3-matrix} gives every E3 probe and control at level 6 together with the level-2 and level-12 references; Table~\ref{tab:e2-dose} gives the E2 solve shares per level, policy and dose; Table~\ref{tab:equivariance} gives the exploratory equivariance result per category of symmetry. All three are descriptive; the confirmatory contrasts are those of Table~\ref{tab:primary}.

\begin{table}[ht]
\centering
\caption{E3 probes and controls on the TEST puzzles (Brier loss of the nine candidate bits; lower is better). The relative difference is taken against the reference row of the same construction. Construction: whole-layer (cell and its peers), per-cell, input (raw puzzle, peer counts, surface features) or label-free.}
\label{tab:e3-matrix}
\scriptsize
\setlength{\tabcolsep}{2pt}
\begin{tabular}{llrlrr}
\toprule
Row & Construction & Brier & Reference & Relative difference (\%) & Parameters \\
\midrule
wl\_L06.generated & whole\_layer & 0.2307 & -- & -- & 5346 \\
wl\_L06.permuted & whole\_layer & 0.2335 & wl\_L06.generated & 1.3 & 5346 \\
wl\_L06\_controls.random & whole\_layer & 0.2337 & wl\_L06.generated & 1.3 & 5346 \\
wl\_L06\_controls.label\_permuted & whole\_layer & 0.2328 & wl\_L06.generated & 1.0 & 5346 \\
wl\_teacher\_L06.teacher & whole\_layer & 0.2307 & wl\_L06.generated & 0.0 & 5346 \\
wl\_teacher\_L12.teacher & whole\_layer & 0.1834 & wl\_L12.generated & 0.0 & 5346 \\
raw\_projected.linear & input & 0.0014 & wl\_L06.generated & \ensuremath{-}99.4 & 2628 \\
raw\_unprojected.linear & input & 0.0105 & wl\_L06.generated & \ensuremath{-}95.4 & 2628 \\
pc\_L06.linear & per\_cell & 0.2310 & -- & -- & 3042 \\
pc\_L06.mlp & per\_cell & 0.2309 & pc\_L06.linear & 0.0 & 22217 \\
controls.raw\_linear & input & 0.0105 & pc\_L06.linear & \ensuremath{-}95.4 & 2628 \\
controls.raw\_mlp & input & 0.0000 & pc\_L06.mlp & \ensuremath{-}100.0 & 19273 \\
controls.peer\_counts\_linear & input & 0.0142 & pc\_L06.linear & \ensuremath{-}93.9 & 819 \\
controls.surface\_linear & input & 0.2253 & pc\_L06.linear & \ensuremath{-}2.5 & 828 \\
controls.permuted\_linear & per\_cell & 0.2326 & pc\_L06.linear & 0.7 & 3042 \\
controls.random\_linear & per\_cell & 0.2330 & pc\_L06.linear & 0.9 & 3042 \\
controls.label\_permuted\_linear & per\_cell & 0.2326 & pc\_L06.linear & 0.7 & 3042 \\
controls.teacher\_L02\_linear & per\_cell & 0.2321 & pc\_L02.linear & 0.0 & 3042 \\
controls.teacher\_L06\_linear & per\_cell & 0.2310 & pc\_L06.linear & 0.0 & 3042 \\
controls.teacher\_L12\_linear & per\_cell & 0.1860 & pc\_L12.linear & 0.1 & 3042 \\
controls.rule & input & 0.0000 & -- & -- & -- \\
controls.marginal & input & 0.2325 & -- & -- & -- \\
\bottomrule
\end{tabular}

\end{table}

\begin{table}[ht]
\centering
\caption{E2 solve shares by level, policy and dose (descriptive). Branch A: clean run; C: replaced codes with replayed later levels; D: replaced codes with regenerated later levels. Level~1 enumerates every alternative value, weighted by the policy (runs: eligible puzzles); the other levels use three intervention seeds per puzzle (runs: puzzle--seed pairs).}
\label{tab:e2-dose}
\footnotesize
\begin{tabular}{rllrrrrr}
\toprule
Level & Policy & Dose & Records & Excluded (\%) & Solve A (\%) & Solve C (\%) & Solve D (\%) \\
\midrule
1 & frequency & enumerated & 600 & 0.0 & 87.7 & 74.9 & 76.5 \\
1 & uniform & enumerated & 600 & 0.0 & 87.7 & 68.0 & 69.9 \\
2 & frequency & 1 & 1800 & 0.0 & 87.7 & 87.7 & 87.7 \\
2 & frequency & 2 & 1800 & 0.0 & 87.7 & 87.7 & 87.3 \\
2 & uniform & 1 & 1800 & 0.0 & 87.7 & 87.3 & 87.1 \\
2 & uniform & 2 & 1800 & 0.0 & 87.7 & 86.8 & 86.6 \\
6 & frequency & 1 & 1800 & 0.0 & 87.7 & 87.7 & 87.6 \\
6 & frequency & 2 & 1800 & 0.0 & 87.7 & 87.7 & 87.0 \\
6 & frequency & 4 & 1800 & 0.0 & 87.7 & 87.7 & 87.3 \\
6 & frequency & 8 & 1800 & 0.0 & 87.7 & 87.7 & 86.1 \\
6 & frequency & 13 & 1800 & 0.0 & 87.7 & 87.7 & 85.6 \\
6 & uniform & 1 & 1800 & 0.0 & 87.7 & 87.7 & 87.3 \\
6 & uniform & 2 & 1800 & 0.0 & 87.7 & 87.7 & 86.9 \\
6 & uniform & 4 & 1800 & 0.0 & 87.7 & 87.7 & 86.6 \\
6 & uniform & 8 & 1800 & 0.0 & 87.7 & 87.7 & 86.9 \\
6 & uniform & 13 & 1800 & 0.0 & 87.7 & 87.7 & 86.0 \\
11 & frequency & 1 & 1800 & 0.0 & 87.7 & 87.6 & 87.7 \\
11 & frequency & 4 & 1800 & 0.0 & 87.7 & 87.6 & 87.4 \\
11 & frequency & 11 & 1800 & 0.0 & 87.7 & 87.4 & 87.3 \\
11 & frequency & 21 & 1800 & 0.0 & 87.7 & 86.7 & 86.5 \\
11 & frequency & 41 & 1800 & 0.0 & 87.7 & 85.4 & 85.9 \\
11 & frequency & 67 & 1800 & 0.0 & 87.7 & 83.6 & 84.2 \\
11 & uniform & 1 & 1800 & 0.0 & 87.7 & 87.7 & 87.6 \\
11 & uniform & 4 & 1800 & 0.0 & 87.7 & 87.6 & 87.5 \\
11 & uniform & 11 & 1800 & 0.0 & 87.7 & 87.5 & 87.3 \\
11 & uniform & 21 & 1800 & 0.0 & 87.7 & 86.8 & 86.9 \\
11 & uniform & 41 & 1800 & 0.0 & 87.7 & 86.3 & 85.9 \\
11 & uniform & 67 & 1800 & 0.0 & 87.7 & 84.4 & 84.9 \\
\bottomrule
\end{tabular}

\end{table}

\begin{table}[ht]
\centering
\caption{Exploratory level-12 equivariance per category of symmetry: mean deviation between the readout of a transformed puzzle and the transformed readout of the original, share of transformed puzzles within the fixed tolerance, and Brier loss on transformed and original puzzles.}
\label{tab:equivariance}
\footnotesize
\begin{tabular}{lrrrr}
\toprule
Category & Mean deviation & Within tolerance (\%) & Brier transformed & Brier original \\
\midrule
all & 0.0791 & 0.0 & 0.1837 & 0.1839 \\
digit & 0.0686 & 0.0 & 0.1834 & 0.1839 \\
geometric & 0.0782 & 0.0 & 0.1838 & 0.1839 \\
\bottomrule
\end{tabular}

\end{table}
